\section{导读：这一季的两个关键词}

本节先建立整期季报的主线。过去几期《全球大模型季报》更强调智能上限、模型进步和 AGI 主线；这一期开始明显转向产品、入口和商业化。广密给出的两个关键词是“分化”和“产品”。分化指模型公司开始向不同能力区间和产品路线拆开；产品指模型智能必须转化成用户体验、流量、品牌心智和商业化闭环。

这期的核心张力在于：一方面，基础模型前三名还在快速进步，持续吃掉 AI 创业者的机会；另一方面，Magic Moment 的产品窗口仍然存在，只是保鲜期越来越短。创业者面对的是既兴奋又绝望的局面：技术每个月都在变强，但巨头下场也越来越快。

\begin{figure}[H]
\centering
\includegraphics[width=0.96\textwidth]{startup-between-giants.png}
\caption{创业者夹在巨头之间：技术每月进步，但模型公司也不断吃掉机会。自制概念图，依据 00:00:05--00:03:54 对谈内容整理。}
\end{figure}

\begin{knowledgebox}{读图：为什么又兴奋又绝望}
兴奋来自技术进步，绝望来自头部模型公司同时做模型和产品。创业者必须抢在巨头产品化之前找到 Magic Window，或者选择巨头难以快速覆盖的垂直场景。
\end{knowledgebox}

\begin{importantbox}{本期核心命题}
大模型行业从“智能上限竞赛”进入“智能 + 产品 + 入口 + 商业化”的复合竞赛。模型能力仍然重要，但产品壁垒、默认入口、品牌心智、垂直整合和 L4 体验开始决定谁能把智能变成真实价值。
\end{importantbox}

\subsection{本章小结}

EP112 是模型季报从“模型智能”转向“产品战略”的一集。它不只是列举公司动态，而是在回答：头部模型公司开始产品化之后，创业者和投资人还能在哪里下注。

